% Vyžaduje balíčky: graphicx, float.
\section{Shrnutí a návaznost na aplikaci}

Analýza ukázala, že data o stavebních objektech v Brně a okolí jsou
pro vizualizaci vhodná, ale vyžadují pečlivé zacházení s neúplnými
a částečně zástupnými hodnotami. Nejdůležitější zjištění jsou následující.

\begin{itemize}
\item Datová sada obsahuje 124\,567 objektů a 278\,171 verzí za období
  09/2015 až 08/2026. Aktuálně platných je 102\,220 objektů, z toho 72\,\%
  v Brně.
\item Struktura zástavby je koncentrovaná: téměř polovinu objektů tvoří
  rodinné domy, 17\,\% garáže a 8\,\% bytové domy. Zástavba je převážně
  nízká (74\,\% objektů má jedno nebo dvě podlaží) a výrazně se liší
  podle využití i polohy: bytové domy jsou převážně vícepodlažní a ve městě,
  v okolí převládají rodinné domy.
\item Vybavení a stáří zástavby se prostorově liší: Brno má starší
  zástavbu (medián 1980) než okolní obce (2001) a lepší vybavení
  kanalizací a plynem. Nová výstavba je v okolních obcích zhruba dvakrát
  intenzivnější než v Brně vzhledem k množství stávající zástavby.
\item Historie je po většinu času klidná, ale obsahuje několik hromadných
  aktualizací (například 01/2024 s 43\,521 změnami) a vlny zániků objektů,
  které odpovídají čištění registru, ne skutečné změně zástavby.
\item Řada atributů je vyplněna jen u dvou třetin objektů a starší data
  dokončení jsou často zástupná. Zaniklé objekty jsou téměř bez údajů.
\end{itemize}

\subsection{Vizualizace v aplikaci}

Z těchto zjištění vychází návrh vizualizací v aplikaci. Přehled uvádí
tabulka~\ref{tab:vizualizace}: pro každou vizualizaci je uvedeno, jaký atribut
zobrazuje a na jakou otázku uživatele odpovídá.

\begin{table}[H]
\centering
\caption{Navržené vizualizace v aplikaci}
\label{tab:vizualizace}
\begin{tabular}{p{3.6cm}|p{3.4cm}|p{6.6cm}}
Vizualizace & Atribut & Otázka uživatele \tabularnewline \hline
Mapa objektů s barvou podle atributu & využití, podlažnost, rok dokončení & Jaká je zástavba na daném místě? \tabularnewline \hline
Hustota v síti při malém přiblížení & počet objektů v buňce & Kde je zástavba nejhustší? \tabularnewline \hline
Časová osa a posuvník data & platnost verze (\texttt{valid\_from}, \texttt{valid\_to}) & Jak vypadalo území v daném měsíci? \tabularnewline \hline
Časová řada změn, nových a zaniklých objektů & počet událostí podle měsíců & Kdy se v oblasti nejvíc měnilo? \tabularnewline \hline
Detail objektu s historií verzí & všechny atributy a jejich změny & Co se u objektu změnilo a kdy? \tabularnewline \hline
Srovnání obcí & podíly využití, vybavení, nová výstavba na 1\,000 objektů & Čím se obec liší od ostatních? \tabularnewline \hline
Rozdělení hodnot s filtrem & podlažnost, zastavěná plocha, počet bytů & Kolik objektů splňuje zvolená kritéria? \tabularnewline \hline
Indikátor úplnosti & podíl vyplněných údajů & Jak je údaj důvěryhodný? \tabularnewline
\end{tabular}
\end{table}

\subsection{Důsledky pro návrh}

Při návrhu aplikace je třeba zohlednit tato zjištění z analýzy:

\begin{itemize}
\item \textbf{Chybějící hodnoty} je nutné zobrazovat explicitně
  (například šedou barvou a položkou „nezjištěno“ ve filtrech), jinak by
  graf s podlažností nebo vybavením podhodnotil třetinu objektů. Totéž platí
  pro srovnání obcí, které se mají počítat z objektů s vyplněnou hodnotou.
\item \textbf{Zástupná data dokončení} (31.\,12., hranice období) nelze
  zobrazovat jako přesný rok. Starší zástavbu je vhodné členit do období
  a po jednotlivých letech zobrazovat až objekty z roku 2001 a novější.
\item \textbf{Hromadné aktualizace} v historii mohou časovou osu zkreslit
  (jeden měsíc má desítky tisíc změn). Řada událostí musí proto mít
  volitelnou logaritmickou osu, případně odlišit hromadné měsíce.
\item \textbf{Zánik objektu} nelze zobrazovat jako zbourání. Objekty, jejichž
  platnost skončila, mají minimum údajů, takže se v aplikaci označují jako
  „vyřazeno z evidence“ a z map se při výchozím zobrazení vynechávají.
\item \textbf{Měřítko mapy}: při velkém území (desítky tisíc bodů) je
  vhodná agregace do šestiúhelníkové sítě, při přiblížení se přechází na
  jednotlivé objekty (u 2,2\,\% objektů bez obrysu na definiční body).
\item \textbf{Mezery ve zdrojích} (chybějící snapshoty, poškozený měsíc Brna
  v 05/2016) se mají v časové ose vyznačit, aby se jejich dopad nepřipisoval
  změnám v zástavbě.
\item \textbf{Rozšíření územím}: analýza se týkala Brna a okolí, ostatní
  regiony má pipeline schopná zpracovat stejným postupem, ale podíly
  vyplněnosti atributů se mohou lišit.
\end{itemize}
